% Part: lambda-calculus
% Chapter: introduction
% Section: free-variables

\documentclass[../../../include/open-logic-section]{subfiles}

\begin{document}

\olfileid{lam}{syn}{fv}
\olsection{自由变元}

λ-演算是关于函数的理论，而函数是通过λ-抽象产生的。
直观地说，$\lambd[x][M]$ 是这样的函数：当自变元被指派给~$x$ 时，其值由~$M$ 给出。
但是，并非~$M$ 中~$x$ 的每一次出现都是相关的：如果 $M$ 包含另一个抽象 $\lambd[x][N]$，
那么 $N$ 中~$x$ 的出现与 $\lambd[x][N]$ 相关，而与 $\lambd[x][M]$ 无关。
因此，$\lambd[x][M]$ 内部的λ-抽象 $\lambd[x]$ \emph{约束}了~$M$ 中那些尚未被另一个λ-抽象约束的~$x$ 的出现——即~$M$ 中~$x$ 的\emph{自由}出现。

\begin{defn}[作用域]
如果 $\lambd[x][M]$ 出现在项~$N$ 内部，则~$N$ 的相应出现即为该~$\lambd[x]$ 的\emph{作用域}。
\end{defn}


\begin{defn}[自由出现与约束出现]
  变元 $x$ 在项~$M$ 中的一次出现是\emph{自由的}，如果它不在任何 $\lambd[x]$ 的作用域内；否则就是\emph{约束的}。
  变元~$x$ 在 $\lambd[x][M]$ 中的一次出现被开头的 $\lambd[x]$ \emph{约束}，当且仅当 $x$ 在~$M$ 中的这次出现是自由的。
\end{defn}

\begin{ex}
  在 $\lambd[x][x y]$ 中，$x$ 和 $y$ 都在 $\lambd[x]$ 的作用域内，因此 $x$ 被 $\lambd[x]$ 约束。
  由于 $y$ 不在任何 $\lambd[y]$ 的作用域内，所以它是自由的。
  在 $\lambd[x][x x]$ 中，$x$ 的两次出现都被~$\lambd[x]$ 约束，因为它们在 $xx$ 中都是自由的。
  在$((\lambd[x][xx])x)$中，最后一次出现的~$x$ 是自由的，因为它不在任何 $\lambd[x]$ 的作用域内。
  在$\lambd[x][(\lambd[x][x])x]$中，第一个 $\lambd[x]$ 的作用域是 $(\lambd[x][x])x$，第二个 $\lambd[x]$ 的作用域是倒数第二次出现的~$x$。
  在 $(\lambd[x][x])x$ 中，最后一次出现的 $x$ 是自由的，倒数第二次出现的是约束的。
  因此，在$\lambd[x][(\lambd[x][x])x]$中，倒数第二次出现的 $x$ 被第二个 $\lambd[x]$ 约束，而最后一次出现的被第一个 $\lambd[x]$ 约束。
\end{ex}

对于一个项 $P$，我们可以检查其中所有变元的出现，从而得到一个自由变元集合。
这个集合记作 $\FV{P}$，其自然定义如下：

\begin{defn}[项的自由变元] \ollabel{def:fv}
  项的自由变元集归纳定义如下：
  \begin{enumerate}
    \item $\FV{x} = \{x\}$ \ollabel{def:fv1} 
    \item $\FV{\lambd[x][N]} = \FV{N} \setminus \{x\}$    \ollabel{def:fv2}
    \item $\FV{PQ} = \FV{P} \cup \FV{Q}$ \ollabel{def:fv3}
  \end{enumerate}
\end{defn}

\begin{prob}
  \begin{enumerate}
  \item 指出此项中 $\lambd[g]$ 和两个 $\lambd[x]$ 的作用域：$\lambd[g][(\lambd[x][g (x x)]) \lambd[x][g (x x)]]$。
  \item 在$\lambd[g][(\lambd[x][g (x x)]) \lambd[x][g (x x)]]$中，所有变元的出现都是被约束的吗？它们分别被哪些抽象约束？
    \item 给出$\FV{\lambd[x][(\lambd[y][(\lambd[z][x y]) z]) y]}$
  \end{enumerate}
\end{prob}

\begin{explain}
自由变元就像是对外部世界（即\emph{环境}）的引用，
一个包含自由变元的项可以看作是一个部分规定的项，
因为它的行为取决于我们如何设置环境。
例如，在项 $\lambd[x][f x]$ 中，它接受一个自变元~$x$ 并返回该自变元的 $f$ 值，
此时变元~$f$ 是自由的。
这个项的值依赖于它所处的环境，特别是该环境中 $f$ 的值。

如果我们对此项再作一次抽象，就得到$\lambd[f][\lambd[x][f
    x]]$。
这个项不再依赖于环境变元 $f$，因为它现在指称一个接受两个自变元并返回将第一个应用于第二个之结果的函数。
改变环境中的 $f$ 不会对这个项的行为产生任何影响，因为这个项只会使用任何作为自变元传入的对象，而不会使用环境中 $f$ 的值。
\end{explain}

\begin{defn}[闭项，组合子]
  没有自由变元的项称为\emph{闭项}，或\emph{组合子}。
\end{defn}

\begin{lem}\ollabel{lem:fv}
  \begin{enumerate}
    \item \ollabel{lem:fv-abs} 若 $y \neq x$，则$y \in
      \FV{\lambd[x][N]}$。
    \item \ollabel{lem:fv-app} $y \in \FV{PQ}$ 当且仅当 $y \in \FV{P}$ 或 $y \in \FV{Q}$。
    \end{enumerate}
\end{lem}

\begin{proof}
  练习。
\end{proof}

\begin{prob}
  证明 \olref[lam][syn][fv]{lem:fv}。
\end{prob}

\end{document}